% Сборка из каталога встречи: latexmk talk.tex
\documentclass[aspectratio=169,11pt]{beamer}
\usepackage{fontspec}
\setsansfont{Arial}
\setmainfont{Arial}
\usepackage[russian]{babel}
\usepackage{tikz,booktabs}
\definecolor{ink}{HTML}{193745}
\definecolor{sea}{HTML}{087B83}
\definecolor{rust}{HTML}{B85726}
\setbeamercolor{normal text}{fg=ink,bg=white}
\setbeamercolor{frametitle}{fg=ink}
\setbeamertemplate{navigation symbols}{}
\setbeamerfont{frametitle}{size=\Large,series=\bfseries}
\setbeamersize{text margin left=8mm,text margin right=8mm}
\setbeamertemplate{footline}{\hfill{\color{sea}\scriptsize\insertframenumber}\hspace{8mm}\vspace{3mm}}
\graphicspath{{figures/}}
\newcommand{\lead}[1]{\textcolor{sea}{\textbf{#1}}}
\begin{document}

\begin{frame}
{\small\color{sea} Рабочая встреча с коллегами из Сколтеха \hfill 25 ноября 2026 г.}
\medskip

{\Large\bfseries Эволюция гауссова импульса в резонаторе\\ и точность её воспроизведения методом PINN}
\medskip

\begin{columns}[c,onlytextwidth]
\begin{column}{0.61\textwidth}
\lead{Уравнение Лугиато--Лефевера}
\[
 \partial_t\psi=F-(1+i\delta)\psi-i\frac{\beta_2}{2}\partial^2_\theta\psi+i|\psi|^2\psi .
\]
На кольце длины $2\pi$ задаём импульс\\
$\psi(\theta,0)=A\exp[-\theta^2/(2\sigma^2)]$.\\[1ex]
$F=2,\quad\delta=3,\quad\beta_2=-0{,}1$.
\end{column}
\begin{column}{0.35\textwidth}
\centering
\begin{tikzpicture}[x=.71cm,y=1.14cm,font=\small]
\coordinate (peak) at (0,1.6);
\coordinate (width) at (1,{1.6*exp(-0.5)});
\draw[->,gray] (-3,0)--(3,0) node[right] {$\theta$};
\draw[->,gray] (0,0)--(0,2.05) node[above] {$|\psi_0|$};
\draw[sea,thick,domain=-3:3,samples=121] plot(\x,{1.6*exp(-\x*\x/2)});
\draw[gray,densely dashed] (-2.3,1.6)--(peak);
\draw[<->,rust] (-2.3,0)--(-2.3,1.6) node[midway,left] {$A$};
\draw[gray,densely dashed] (width)--(1,0);
\draw[<->,rust] (0,-.18)--(1,-.18) node[midway,below] {$\sigma$};
\fill[sea] (width) circle (1.5pt) node[above right,ink] {$A e^{-1/2}$};
\end{tikzpicture}
\end{column}
\end{columns}
\medskip
Как амплитуда и ширина определяют исход? Насколько точно нейросеть воспроизводит формирование импульса?
\end{frame}

\begin{frame}{В 28 из 64 расчётов сформировался один пик}
\centering\includegraphics[width=.94\linewidth]{outcomes.pdf}\par
\raggedright\small
На окне $80\le t\le100$ в 31 случае поле однородно, в 28 стационарно с одним пиком. Кольцами отмечены начальные условия для профилей справа.
\end{frame}

\begin{frame}{При накачке $F=2{,}2$ поле ещё меняется во всех 64 узлах}
\small Фиксируем отстройку $\delta=3$. Для каждого значения накачки повторяем расчёт на той же сетке начальных импульсов и сравниваем поле на окне $80\le t\le100$.
\medskip
\centering\includegraphics[width=\linewidth]{pump-series.pdf}
\end{frame}

\begin{frame}{При отстройке $\delta=3{,}5$ один пик сохраняется в 15 узлах}
\small Фиксируем накачку $F=2$. При $\delta=2{,}5$ все 64 решения ещё меняются на окне $80\le t\le100$. При $\delta=3{,}5$ все 64 удовлетворяют критерию стационарности.
\medskip
\centering\includegraphics[width=\linewidth]{detuning-series.pdf}
\end{frame}

\begin{frame}{Три из пяти переходных режимов пришли к одному пику}
\small Пять начальных импульсов, для которых поле ещё менялось при $T=100$,
прослежены до $T=200$ или $400$ при исходных $F=2$, $\delta=3$.\par\medskip
\centering
\begin{tabular}{rrrl}
\toprule
$A$ & $\sigma$ & $T$ & Состояние на последних 20 единицах времени\\
\midrule
2,4 & 0,480 & 200 & Один стационарный пик\\
2,8 & 0,700 & 400 & Один стационарный пик\\
3,2 & 0,700 & 400 & Один стационарный пик\\
\addlinespace
2,8 & 0,480 & 400 & Два пика, расстояние между ними меняется\\
3,2 & 0,329 & 400 & Два пика, расстояние между ними меняется\\
\bottomrule
\end{tabular}
\end{frame}

\begin{frame}{Как меняется расстояние между пиками при $A=2{,}8$?}
\centering\includegraphics[width=.94\linewidth,height=.80\textheight,keepaspectratio]{dynamics-a28.pdf}
\end{frame}

\begin{frame}{Слабые волны расходятся от центрального импульса}
\centering\includegraphics[width=.94\linewidth]{waves.pdf}\par
\raggedright\small
Спектральный расчёт, $A=1{,}6$, $\sigma=0{,}329$.
Справа показана $\partial_t(\arg\psi-\arg\psi_{\rm bg})$ в рад/ед.~$t$.
Фон $\psi_{\rm bg}(t)$ получен при той же накачке из нулевого однородного поля.
\end{frame}

\begin{frame}{До $T=5$ максимальная ошибка поля составила 1,17\%}
\centering\includegraphics[width=.97\linewidth]{pinn.pdf}\par
\raggedright\small
Одна PINN на $[0,5]$, $A=1{,}6$, $\sigma=0{,}329$. Обучение: 131\,072 точки невязки, 2 млн шагов Adam. Ошибка вычислена по комплексному полю относительно SSFM.
\end{frame}

\begin{frame}{Как точнее решить задачу с тем же числом точек?}
Расчёт PINN до $T=5$ занял \lead{17 ч 15 мин} на RTX 4090 при 131\,072 точках невязки и 2 млн шагов. Можно ли лучше выбрать точки, в которых сеть проверяет уравнение?
\medskip

\begin{columns}[c,onlytextwidth]
\begin{column}{.53\textwidth}
\centering\includegraphics[width=\linewidth]{rad-peaks.pdf}
\end{column}
\begin{column}{.44\textwidth}
\lead{Пример при $T=4$}\\[.4ex]
\small В обоих запусках 8\,192 точки и 500 тыс. шагов. Каждые 100 шагов обновляем выборку. При RAD чаще выбираем точки с большой невязкой.\\[1ex]
Максимальная ошибка поля при равномерной выборке составила 73,9\% за 71 мин, при RAD --- 2,8\% за 78 мин.
\end{column}
\end{columns}
\end{frame}
\end{document}
